\chapter{Introduction}
\label{chap:introduction}

Populism presents a genuine empirical puzzle for the study of corruption. Populist parties define
politics as a moral conflict between a virtuous people and a self-serving elite
\citep{mudde2004}. This framing makes anti-corruption rhetoric electorally natural: established
parties, bureaucracies, courts, and media can all be portrayed as parts of a captured system. Yet the
same majoritarian logic may weaken horizontal accountability once populists enter government.
Concentrating executive authority, politicising appointments, and discrediting independent checks can
create opportunities for institutional capture \citep{chesterley2018,batory2016}. The relationship can
also run in reverse. Persistent corruption and declining trust may increase demand for outsiders who
promise to displace established elites \citep{guriev2022,kubbe2020}.

These mechanisms imply that a simple cross-country correlation is not enough. Countries in which
populists govern differ systematically in economic development, democratic institutions, regional
history, and baseline corruption. Moreover, both governing-party ideology and corruption change
slowly. A positive pooled relationship can therefore reflect stable differences between countries,
while a same-year fixed-effects estimate cannot establish which variable came first. Finally, an
aggregate corruption index may hide divergent effects across the executive, legislature, judiciary,
and public administration. Existing research likewise emphasises that populism is heterogeneous and
its institutional effects need not be identical across ideological or political settings
\citep{huber2017,engler2024}.

This report addresses these problems through three research questions:

\begin{enumerate}[label=\textbf{RQ\arabic*.}]
  \item Within a country, are years of more populist government associated with higher political
  corruption after accounting for country baselines, common year shocks, income, and democracy?
  \item Does governing-party populism predict subsequent corruption, does corruption predict
  subsequent populism, or is there no robust temporal ordering?
  \item Do the results differ across executive, public-sector, legislative, and judicial corruption?
\end{enumerate}

The empirical contribution is not a search for a model that produces the strongest result. Instead,
each method answers a distinct question. Pooled and between-country estimates describe cross-national
differences. Two-way fixed effects isolate within-country co-movement. Granger-style models ask whether
one variable adds predictive information beyond the other's own persistence. Event studies, placebo
leads, first differences, alternative electoral samples, and multiple-testing correction delimit the
claims that can be sustained.

The findings are deliberately nuanced. The pooled populism coefficient is positive, but statistically
imprecise once country dependence is recognised. Within countries, the coefficient becomes small and
negative. It is significant in the controlled static model, but its significance is sensitive to
alternative measurement and differencing choices. The composite index provides no robust evidence for
either temporal direction. Decomposition clarifies that judicial corruption accounts for the strongest
negative same-year pattern, whereas a small positive two-year legislative signal is the only
dimension-specific temporal exception. That signal survives the primary Granger and false-discovery
tests but disappears in first-difference and election-only designs.

The central conclusion is therefore not that populism reduces corruption, nor that corruption is
irrelevant. Rather, the headline cross-country association does not translate into a general
within-country or temporal effect. The evidence is more consistent with strong contextual confounding,
slow-moving measurement, and limited institution-specific dynamics. This negative result is
substantively useful: it identifies where a widely assumed relationship does and does not survive a
more demanding panel design.

The remainder of the report proceeds as follows. \Cref{chap:data} describes the data, measurement,
sample construction, and exploratory evidence. \Cref{chap:methods} sets out the estimands and model
specifications. \Cref{chap:results} presents the association, temporal, and decomposition findings,
then discusses limitations and implications. The appendix documents robustness and reproducibility.
